\section{DỮ LIỆU VÀ THIẾT LẬP THỰC NGHIỆM}
\label{sec:experimental}
% =============================================================================

\subsection{Phần cứng \& Phần mềm}

Tất cả các thử nghiệm được thực hiện trên một máy trạm có CPU Intel Core i7-12700, 32~GB RAM, và GPU NVIDIA RTX~3060 12~GB.
Môi trường phát triển sử dụng Python~3.12 với các thư viện chính:
Pandas~2.1 và NumPy~1.26 cho xử lý dữ liệu,
Scikit-learn~1.4 cho chuẩn hóa và chia tách,
PyWavelets~1.6 cho biến đổi wavelet,
TensorFlow~2.21 cho huấn luyện GAN,
và Matplotlib~3.8 cùng Seaborn~0.13 cho trực quan hóa.
Mã nguồn và pipeline được quản lý bằng Git với cấu trúc modular (OOP),
đảm bảo khả năng tái sử dụng và mở rộng.

\subsection{Chia tách Tập dữ liệu}

Tập dữ liệu gồm 35.040 mẫu ghi nhận liên tục trong năm 2018.
Đối với dữ liệu chuỗi thời gian, chia ngẫu nhiên sẽ phá vỡ cấu trúc phụ thuộc thời gian và gây data leakage.
Do đó, chúng tôi áp dụng \textbf{chia tách theo thời gian} (temporal split):
80\% đầu (tháng 1--9) dùng cho huấn luyện và 20\% cuối (tháng 10--12) dùng cho kiểm tra.
Đối với mô-đun GAN, chỉ các mẫu bất thường (minority, 2.388 mẫu) được dùng để huấn luyện generator
nhằm học phân phối của lớp thiểu số và sinh thêm mẫu synthetic;
toàn bộ mẫu bình thường được giữ nguyên làm tập đa số.
Chiến lược này đảm bảo GAN không bị chi phối bởi phân phối của lớp đa số.

\subsection{Kiểm tra Tính toàn vẹn Dữ liệu}

Tập dữ liệu gốc 35.040 mẫu và dữ liệu thờ tiết 8.760 mẫu đều không chứa giá trị thiếu (zero missing).
Điều này phản ánh chất lượng thu thập cao từ hệ thống SCADA và API khí tượng Open-Meteo.

Dù vậy, chúng tôi vẫn triển khai module phân tích cơ chế thiếu dữ liệu (MCAR/MAR/MNAR)
như một công cụ audit phòng ngừa trong pipeline.
Khi dữ liệu mở rộng trong tương lai hoặc tích hợp thêm cảm biến phụ trợ,
module này sẽ tự động phân loại cơ chế missing và đề xuất phương pháp xử lý phù hợp
(mục~\ref{sec:methodology}).

\subsection{Kiểm tra rò rỉ dữ liệu}

Một cuộc audit toàn diện được thực hiện để đảm bảo \textit{zero data leakage} trong pipeline tiền xử lý.
Các điểm kiểm tra bao gồm:

\begin{enumerate}
    \item \textbf{Rolling window:} Xác nhận \texttt{center=False} trong mọi hàm tính rolling statistics. Việc sử dụng \texttt{center=True} sẽ đưa thông tin tương lai vào giá trị hiện tại, gây leakage nghiêm trọng.
    \item \textbf{DWT window:} Kiểm tra cửa sổ wavelet chỉ bao gồm $[t-w+1, t]$ (quá khứ + hiện tại), không có look-ahead.
    \item \textbf{Lag features:} Xác nhận chỉ sử dụng \texttt{shift(lag)} với $\text{lag} > 0$ (nhìn về quá khứ).
    \item \textbf{Weather merge:} Xác nhận không sử dụng \texttt{bfill} (back-fill) --- chỉ dùng dữ liệu thời tiết đã quan sát tại hoặc trước timestamp.
    \item \textbf{GAN scaler:} \texttt{MinMaxScaler.fit\_transform()} được fit trên toàn bộ dữ liệu trước khi huấn luyện GAN. Điều này \textbf{không} gây leakage cho mô hình downstream vì: (i) downstream model sẽ fit scaler riêng trên tập train; (ii) scaler là phép biến đổi tuyến tính đơn giản; (iii) GAN chỉ sinh synthetic data, không phải model dự đoán.
    \item \textbf{Anomaly labels:} Rolling mean trong \texttt{label\_leakage} tính trên toàn bộ chuỗi. Đây là việc xác định ground truth dựa trên vật lý (ISO 50001 baseline), không phải feature dùng để huấn luyện model.
\end{enumerate}

Kết quả audit xác nhận pipeline không chứa data leakage.
Mọi transformation đều chỉ sử dụng thông tin quá khứ hoặc hiện tại,
tuân thủ nguyên tắc causal learning cho chuỗi thời gian.

\subsection{Chỉ số Đánh giá}

Hiệu quả của pipeline được đánh giá qua ba nhóm chỉ số:
\begin{itemize}
    \item \textbf{Độ trung thực thống kê:} So sánh trung bình, độ lệch chuẩn, và các phân vị (25\%, 50\%, 75\%) giữa dữ liệu thực và tổng hợp.
    Sai số tương đối dưới 5\% được coi là chấp nhận được.
    \item \textbf{Bảo toàn cấu trúc phân phối:} Sử dụng PCA và t-SNE để chiếu dữ liệu xuống 2D;
độ chồng chéo (overlap) trực quan giữa các điểm thực và tổng hợp phản ánh mức độ bảo toàn phân phối đa chiều.
    \item \textbf{Bảo toàn ma trận tương quan:} Tính sai số Frobenius $\| \mathbf{R}_{\text{real}} - \mathbf{R}_{\text{syn}} \|_F$
giữa ma trận tương quan của dữ liệu thực và synthetic.
Giá trị nhỏ chứng tỏ các mối quan hệ tuyến tính giữa đặc trưng được bảo toàn tốt.
\end{itemize}

Ngoài ra, chúng tôi sử dụng 49 unit tests (pytest) để đảm bảo tính đúng đắn của các module core,
bao gồm data loader, assertions, và anomaly labeling.

% =============================================================================
